\begin{abstract}
Constructing and serving knowledge graphs (KGs)
is an iterative
and human-centered process involving
on-demand programming and analysis.
In this paper, we present \system, a programmable and interactive widget 
library that facilitates human-in-the-loop knowledge acquisition and integration to enable
continuous curation a knowledge graph (KG). 
\system provides a seamless environment within computational notebooks where
data scientists explore a KG to identify opportunities for acquiring new knowledge and 
verify recommendations provided by AI agents for integrating the acquired knowledge in the KG.
We refined \system through participatory
design and conducted case studies
in a real-world setting for evaluation.
The case-studies show that introduction of \system
within an existing HR knowledge graph construction and serving platform
improved the user experience of the experts and helped eradicate inefficiencies
related to knowledge acquisition and integration tasks.


\end{abstract}

%submitted
% 


%Alternate:
% Data science workflows are human-centered processes involving on-demand programming and analysis. While programmable and interactive interfaces such as widgets embedded within computational notebooks are suitable for these workflows, they lack robust state management capabilities and do not support user-defined customization of the interactive components. The absence of such capabilities hinders the reusability and reproducibility of workflows while limiting the scope of exploration of the end-users. In response, we developed \system, a framework for authoring interactive widgets within computational notebooks that enables reusable, reproducible, and customizable data science workflows. The framework redesigns widgets to support fine-grained interaction history management, reusable states, and user-defined customizations. We conducted three case studies in a real-world knowledge graph construction and serving platform to evaluate the effectiveness of these widgets. Based on the observations, we discuss future implications of employing \system widgets for general-purpose data science workflows. 

 